\subsection{Object Copies Are Not Independent Computations}
\label{sec:copy}

The original \filepath{lib/copy.abclr2} creates another object from an
existing object's scripts, state information, and pending messages.
For an OOP reader, the important question is not just ``is the identity
different?'' but ``which references and pending effects remain shared?''
This library also underlies \code{report} and can move an object into
the \code{watching} group defined by \code{watch}.

\paragraph{What is shared.}
Copying preserves local state values but is neither deep nor atomic: new local bindings retain the
same referenced data, while outer lexical environments, scripts, and queued
reply destinations remain shared. The examples below separate these effects;
do not interpret them as independent replay of an active computation.

\subsubsection{Read the Original Fibonacci Experiment}

The library's final block comment defines a recursive Fibonacci object.
It creates a collector for each nontrivial request and sends two requests
to itself: both target \code{Me}, not newly created worker objects.
The later Pcontrol project (Section~\ref{sec:pcontrol-fib}) uses newly
created children instead; you do not need that project to follow this one.

\wholelisting[firstline=164,lastline=179]{lib/copy.abclr2}

The collector captures \code{R}, the reply destination of the original
request. Its local \code{half-result} stores the first answer; the next
answer is added to it and sent to \code{R}. The base cases are 0 and 1.
Without copying an active computation, \code{[fib <== 5]} returns 5.
The library then proposes this experiment:

\wholelisting[firstline=181,lastline=187]{lib/copy.abclr2}

To replay it, load \code{lib/copy} with \code{inc-lib}, enter the
original \code{fib} definition, then enter this expression. The automated
\code{copy-check} target extracts both source ranges without editing them.
The \code{dotimes} loop is a scheduling opportunity, not a timed delay.
The copy can inherit requests already waiting for \code{fib}; those
requests still carry references to the original reply continuations and
collectors. Starting the copy can therefore deliver additional answers
into the original computation. The original comment warns that the result
can seem ``funny.''

Do not expect one fresh, independent Fibonacci answer from this session.
In the checked run the same \code{Fib(5)=} label appeared several times
with values including 4, 5, and 6. Those values are observations, not a
specified sequence to reproduce on every schedule. A collector designed
for two answers cannot distinguish extra answers caused by duplicated
work. A successful ordinary \code{[fib <== 5]} after the experiment does
not certify independent continuation copying.

\subsubsection{Make Sharing Visible with a Small Object}

The supplemental \filepath{tutorial/examples/11-copy.abclr2} puts a
counter binding in an \emph{outer lexical environment}, not in an object's
local \code{state} frame:

\wholelisting[firstline=4,lastline=10]{tutorial/examples/11-copy.abclr2}

The outer assignment retains the object in the top-level binding
\code{copy-source}. The name \code{copy-subject} inside the \code{let}
is local to that environment; a printed object name is not a guarantee
that a top-level variable of that name exists.
In a fresh session, try:

\begin{codeblock}
\begin{lstlisting}
(inc "tutorial/examples/11-copy" :use-source t)
[clone := [copier <== [:copy copy-source]]]
(eq clone copy-source)                 ; => NIL
[copy-source <== :increment]           ; => 1
[clone <== :value]                     ; => 1
[clone <== :increment]                 ; => 2
[copy-source <== :value]               ; => 2
\end{lstlisting}
\end{codeblock}

The objects and their metaobjects have distinct identities, but the
scripts are reused and the outer environment is retained. Both objects
therefore update the same binding \code{shared}. This is not a deep copy
of all reachable Lisp data or of captured application objects.
It is also not the proposed \code{COPY} environment semantics discussed
as future work in the historical guide.

\subsubsection{Follow the Copier's Construction Protocol}

Start with the \code{:group} handler, which the default copy request
also uses:

\wholelisting[firstline=44,lastline=54]{lib/copy.abclr2}

It first queries the original metaobject for its queue, state environment,
and scripts. It asks that environment for its local association list and
outer \code{:lexEnv}. Then \code{alist->stateSpec} builds an object
creation specification, and the selected group manager receives
\code{[:new stateSpec lexEnv scripts]}. The reply is a \emph{metaobject}.
The continuation gives the new object a display name, arranges its pending
messages through \code{message-setter}, and forwards the final object
to the copier's caller.

\noindent\begin{minipage}{\linewidth}
\begin{longtable}{L{0.30\textwidth}L{0.60\textwidth}}
\textbf{Part} & \textbf{Copy protocol}\\\hline
Object/metaobject identity & Creates new identities; renaming does not redirect
existing application references.\\
Scripts & Passes the existing script objects to the new construction.\\
Local state & Reconstructs an initializer specification from current values;
creates separate bindings but retains referenced data.\\
Outer lexical environment & Reuses the original environment reference.\\
Pending messages & Enqueues the original message/reply/sender records on the
copy; leaves the original queue in place.\\
Active execution & Does not clone an executing evaluator context or
recursively replace captured continuations and collectors.\\
\end{longtable}
\end{minipage}

These are several reflective requests, not an atomic snapshot protocol.
If an application continues to change while they run, do not assume the
collected queue and state describe one indivisible instant.
Custom metaobjects must support the queried reflection protocol.

\paragraph{How local-state values are preserved.}
For a binding \code{(n . 7)}, the helper captures 7 in a zero-argument
function and uses that function as the operand of \code{:self-evaluate}.
The primary evaluator calls it with \code{funcall} to obtain the value.
This also preserves a symbol or function as data rather than accidentally
calling it. The binding is new, not the referenced value: assigning the
copy's numeric slot leaves the source's slot unchanged; mutating a shared
list is visible through both. Rebinding that list slot does not rebind its
aliases. Initializer order and the outer environment are unchanged.

\subsubsection{Choose the Group and Evaluator Deliberately}

The library accepts these ordered message patterns. In the table,
\code{obj}, \code{g}, \code{mg}, and \code{ev} stand for evaluated objects:

\noindent\begin{minipage}{\linewidth}
\begin{longtable}{L{0.48\textwidth}L{0.42\textwidth}}
\textbf{Request to \code{copier}} & \textbf{Construction policy}\\\hline
\code{[:copy obj]} & Original group, using that group's current generator
and evaluator.\\
\code{[:copy obj :group g]} & Selected group's generator and evaluator.\\
\code{[:copy obj :meta-gen mg]} & Selected generator, original object's
current evaluator, original group.\\
\code{[:copy obj :meta-gen mg :eval ev]} & Selected generator and evaluator,
original group.\\
\end{longtable}
\end{minipage}

In particular, the simplest copy does \emph{not} necessarily retain an
object-specific evaluator override: it goes through the group's current
creation policy. The explicit \code{:meta-gen} form retrieves the source
object's evaluator instead. The tests distinguish these cases by comparing
evaluator identities. These are message patterns, not unrestricted Common
Lisp keyword arguments; retain the indicated order and combinations.

The macro \code{(copy name)} assigns the copy back to the binding
\code{name}. It also accepts \code{:group} or \code{:meta-gen}, but not both,
and has no \code{:eval} keyword of its own. For example, if
\code{[old-reference := copy-source]} precedes \code{(copy copy-source)},
\code{old-reference} still denotes the original object. Rebinding one name
does not update all references held by other objects.

\subsubsection{Copy Pending Work without Silently Starting It}

The original queue installer is short enough to read completely:

\wholelisting[firstline=150,lastline=157]{lib/copy.abclr2}

It sends \code{:enqueue} for each saved record and replies with
\code{[den newMeta]}, the new application object. The \code{:begin}
send is deliberately commented out. Copying pending work therefore does
not start that work; it also does not remove the original requests.
When a copy already has a nonempty queue, send \code{:begin} to its
metaobject to start it. Do not send \code{:begin} blindly to an empty queue.
An empty-queue copy can receive ordinary application requests normally.
An ordinary query is not a passive queue probe: it can start execution.
Inspect the metaobject's \code{:queue} instead when checking pending work.

The bounded queue exercise stages one pending record explicitly, so it
does not depend on interrupting a running computation:

\begin{codeblock}
\begin{lstlisting}
[object copy-inbox
 (state [values := nil])
 (script
  (=> :values !(reverse values))
  (=> value [values := (cons value values)]))]
[[meta copy-source]
 <= [:enqueue [[:echo 42] copy-inbox copy-source]]]
[queued-clone := [copier <== [:copy copy-source]]]
[copy-inbox <== :values]                ; => NIL
[[meta queued-clone] <= :begin]
[copy-inbox <== :values]                ; => (42)
[[meta copy-source] <= :begin]
[copy-inbox <== :values]                ; => (42 42)
\end{lstlisting}
\end{codeblock}

Here \code{[:enqueue ...]} takes the internal triple
\code{[message reply sender]}, not an ordinary application payload.
Both copies retain \code{copy-inbox} as the reply destination. The
original queue still held its request after the clone ran; starting it
produces the second 42. Enqueueing this way is a controlled reflective
exercise, not the recommended way to submit normal application work.
Do not interpret copying as exactly-once processing, migration that
deletes the source, or independent replay of an active computation.

\paragraph{Evidence and exercise.}
Run \code{make -C tutorial copy-check}. Fifty-two successful probes
cover the original Fibonacci session, identities, shared lexical bindings,
script reuse, queue copying/startup, alias rebinding, and group/evaluator
selection. Twenty-three state checks cover values, reference identity,
independent bindings, shared mutation, and both construction routes. Source extracts,
copies, inputs, and transcripts stay in \filepath{build/copy-tests/}.

Why do both copies increment \code{shared}? They retain the same outer
environment. Why does the inbox receive two replies? Both queues contain
a request addressed to it; only the execution start is separate. As an
exercise, keep the original and a named copy, then change which binding
you query. Predict identity and lexical effects separately. Do not infer
deep-copy semantics from separate object identities.
